\scanpage{42}
\mspage{3}

\[
k(C)\handwritten{\in}H^{\handwritten{3}}(\handwritten{\pi}_0(C),\handwritten{\pi}_1(C)).
\]
En précisant ces réflexions, on trouve que la classification, à $\otimes$-équivalence près, des (Gr)-catégories, se fait précisément en termes d\deleted{\unsure{e}}'un groupe $\handwritten{\pi}_0$, d'un groupe commutatif $\handwritten{\pi}_1$ sur lequel $\handwritten{\pi}_0$ opère, et d'une classe de cohomologie (qui peut être prise arbitraire) dans $H^{\handwritten{3}}(\handwritten{\pi}_0,\handwritten{\pi}_1)$. La loi de \deleted{\unsure{compo}} groupe du $H^{\handwritten{3}}$ admet d'ailleurs une interprétation ``géométrique'' à la Baer, en termes d'opérations sur les (Gr)-catégories.

\uline{Cas particuliers} Dans l'exemple 1, on \inserted{doit} retrouver bien sûr le premier invariant de Postnikoff $\handwritten{\in}H^{\handwritten{3}}(\handwritten{\pi}_1(X),\handwritten{\pi}_2(X))$, avec une interprétation ``géométrique'' nouvelle. Dans l'exemple 2, on trouve $\handwritten{\pi}_0=H^1(X,F)$, $\handwritten{\pi}_1=H^0(X,F)$\handwritten{,} et l'élément du $H^{\handwritten{3}}$ est nul (du fait qu'on a une ``catégorie de Picard stricte'', cf\inserted{\handwritten{Deligne}}\handwritten{[~]}). Dans \deleted{\unsure{les}} le dernier exemple, on a $\handwritten{\pi}_0={}$ groupe des classes d'isomorphisme d'autoéquivalences de $A$, $\handwritten{\pi}_1={}$ groupe des automorphismes du foncteur identique de $A$, et l'invariant $k(C)$ semble nouveau (même dans le cas où $C$ est la catégorie des biModules inversibles sous un Anneau $\underline O_X$, où $\handwritten{\pi}_0$ peut donc s'interpréter comme un ``groupoide de Brandt'' et $\handwritten{\pi}_1$ comme le groupe des ``unités centrales'' de $\underline O_X$). Signalons cependant que lorsque $A$ est la catégorie des torseurs sous un groupe (ordinaire) $F$, on trouve \deleted{un élément de $H^{\handwritten{3}}$\unsure{$($}} $\handwritten{\pi}_0={}$ groupe des automorphismes extérieurs de $F$, $\handwritten{\pi}_1={}$ centre de $F$, et l'élément $k(C)$ du $H^{\handwritten{3}}$ n'est autre que le classique invariant de Mac-Lane. Il resterait à \deleted{établir}\inserted{développer} un lien \inserted{géométrique} \uline{direct} entre la théorie de\deleted{s} Mac-Lane des extensions de groupes non commutatifs\deleted{\unsure{,}} et des liens (``kernels''\deleted{\unsure{)}} dans sa terminologie) auxquels ils sont associés, et la théorie \deleted{\unsure{développée ici|développées ici,}} esquissée ici des (Gr)-catégories et de leurs splittages.

2. \uline{Catégories de Picard}.

On appelle ``catégorie de Picard'' une $\otimes$-catégorie ACU (associative-commutative-unitaire , cf. \handwritten{[~]}) qui est une (Gr)-catégorie pour sa \deleted{\unsure{loi $\otimes$}} $\otimes$-structure AU. On dit que c'est une catégorie de Picard \uline{stricte} si pour tout $X\handwritten{\in}\operatorname{Ob}C$, la symétrie de $X\otimes X$ est l'identité. Ces \deleted{catég} dernières catégories sont étudiées systématiquement (dans le contexte plus général des ``champs de Picards'') par Deligne \handwritten{[~]}.

Soit $C$ une catégorie de Picard. En tant que (Gr)-catégorie, $C$ possède des invariants $\handwritten{\pi}_0$ et $\handwritten{\pi}_1$, mais dans le cas actuel $\handwritten{\pi}_0$ est aussi commutatif et agit trivialement sur $\handwritten{\pi}_{\unsure{0|1}}$. On fera attention ici qu'une équivalence de (Gr)-catégories entre deux catégories de Picard $C,C'$ n'est pas nécessairement compatible avec les contraintes de commutativité, donc n'est pas nécessairement une équivalence de catégories de
% archive sheet number: 41
